\chapter{Mechanical Design}
\label{ch:mechanical}

\section{Overview}

All mechanical parts of the Cubilox Mk1 were designed by the author in Fusion~360 and printed in
standard PLA on a Bambu Lab A1 printer, using the default settings of the Bambu Studio slicer.
The frame of the cube consists of three motor walls carrying the motors and reaction wheels, three
side walls, and eight corner mounts that join the walls with M3 screws and nuts. The electronics
sit on a base plate inside the cube. \Cref{tab:printed_parts} lists all printed parts and
\cref{tab:mech_params} the key parameters.

\begin{table}[htbp]
  \centering
  \caption{3D-printed parts of the Cubilox Mk1.}
  \label{tab:printed_parts}
  \begin{tabular}{lcl}
    \toprule
    Part & Quantity & Function \\
    \midrule
    \cmd{motor\_wall}            & 3 & frame face carrying one motor and its reaction wheel \\
    \cmd{side\_wall}             & 3 & frame face without motor \\
    \cmd{corner\_mount}          & 8 & joins the walls at the corners (M3 nuts press-fit) \\
    \cmd{motor\_mount}           & 6 & holds a motor (fixed with M4 screws) \\
    \cmd{flywheel}               & 3 & reaction wheel \\
    \cmd{sensor\_mount}          & 1 & holds the IMU near the centre of the cube \\
    \cmd{electronic\_base\_plate} & 1 & carries the electronics \\
    \cmd{intermediate\_plate}    & 1 & encloses the battery and carries the circuit board \\
    \bottomrule
  \end{tabular}
\end{table}

\begin{table}[htbp]
  \centering
  \caption{Key mechanical parameters.}
  \label{tab:mech_params}
  \begin{tabular}{ll}
    \toprule
    Parameter & Value \\
    \midrule
    Total mass of the cube                         & \SI{1036}{g} \\
    Mass of one reaction wheel incl.\ screws and nuts & \SI{76.4}{g} \\
    Mass of the battery                            & approx.\ \SI{80}{g} \\
    PLA used for all printed parts                 & \SI{283}{g} \\
    Total print time                               & approx.\ \SI{14}{h} \\
    Edge length of the cube                        & \SI{157}{mm} \\
    \bottomrule
  \end{tabular}
\end{table}

\section{Design for 3D printing}

All parts were designed to be printable without or with very few supports. Overhangs are kept at
printable angles, and horizontal holes and nut pockets use teardrop or pointed shapes, so that the
upper side of the opening is built as a self-supporting overhang instead of a flat bridge.
\Cref{fig:print_optimization} shows two examples.

\begin{figure}[htbp]
  \centering
  \begin{subfigure}[b]{0.42\textwidth}
    \centering
    \includegraphics[width=\textwidth]{print_optimization_example_1.jpg}
    \caption{Nut pocket with a pointed roof, printable without supports.}
    \label{fig:print_opt_1}
  \end{subfigure}\hfill
  \begin{subfigure}[b]{0.42\textwidth}
    \centering
    \includegraphics[width=\textwidth]{print_optimization_example_2.jpg}
    \caption{Detail of a printed frame part with a screw hole and cut-out.}
    \label{fig:print_opt_2}
  \end{subfigure}
  \caption{Examples of geometry optimised for FDM printing.}
  \label{fig:print_optimization}
\end{figure}

\section{Reaction wheels}
\label{sec:wheels}

Each reaction wheel (\cref{fig:wheel}) is a printed ring with four spokes. Six M6 screws, each with
two M6 nuts, are mounted on the rim. They add mass at the largest possible radius and thereby
increase the moment of inertia of the wheel without making the wheel much heavier. One complete
wheel including all screws and nuts weighs \SI{76.4}{g}. The hub is clamped onto the motor shaft
with two screws.

\begin{figure}[htbp]
  \centering
  \includegraphics[width=0.6\textwidth]{reaction_wheel_with_bolts.jpg}
  \caption{Reaction wheel with the M6 screws and nuts on its rim and the two clamping screws at the
    hub.}
  \label{fig:wheel}
\end{figure}

The wheel must sit exactly perpendicular on the motor shaft. A wheel that wobbles causes
vibrations, which disturb the IMU, and it can rub against the frame. During development, a wheel
that rubbed against the frame because of a loose screw produced cracking noises and shocks that
shortened the balancing time considerably (\cref{sec:history}).

\section{Sensor mount and sensor position}

The IMU is mounted as close as possible to the geometric centre of the cube, with each of its
three axes aligned with the axis of one motor. This is why the motors are called X, Y and Z: motor X
rotates the cube about the sensor's X axis, and so on. Placing the sensor near the centre keeps the
accelerations it measures during rotations small.

The sensor mount (\cref{fig:sensor_mount}) was created with the generative design function of
Fusion~360. Despite its organic shape, it can be printed with remarkably little support material.

\begin{figure}[htbp]
  \centering
  \begin{subfigure}[b]{0.36\textwidth}
    \centering
    \includegraphics[width=\textwidth]{mpu_with_holder.jpg}
    \caption{Printed sensor mount with the IMU board.}
    \label{fig:mpu_holder}
  \end{subfigure}\hfill
  \begin{subfigure}[b]{0.56\textwidth}
    \centering
    \includegraphics[width=\textwidth]{sensor_position_cad.png}
    \caption{Section view of the CAD model showing the sensor position.}
    \label{fig:sensor_cad}
  \end{subfigure}
  \caption{Sensor mount and position of the IMU inside the cube.}
  \label{fig:sensor_mount}
\end{figure}

\section{Battery compartment}

The battery sits in a compartment that is closed by the intermediate plate. The circuit board is
screwed on top of this plate with some clearance, so that the cables underneath the board are not
crushed. The compartment is padded with cardboard and folded paper. The padding
fixes the battery in place and absorbs some of the shock when the cube falls over. The
accelerations during a normal fall are usually not high enough to damage the battery, but
securing a LiPo battery against shocks is good practice.
